\documentclass[10pt]{article}
\usepackage{icmlko}

\icmlmeta{IP 파운데이션 월드모델}{50}{이력이 통하는 시장과 통하지 않는 시장}{창작자 7만 6천 건을 색인해 펀딩 매장 노출도를 채웠는데 신호가 정확히 영이었다}

\begin{document}

\twocolumn[
\icmlheader

\begin{abstract}
노트 40부터 미결이던 것을 처리했다. 펀딩의 매장 노출도를 창작자 사전 프로젝트
수로 잡았는데 표본 안에서 세니 평균 0.01로 사실상 상수였다 --- 창작자 7만 명이
대부분 1건이라 표본 400건 안에서는 중복이 거의 없다. 목록 3,876쪽을 전부 훑어
\textbf{창작자 39,587명 · 프로젝트 76,045건}을 색인했고, 2건 이상인 창작자가
10,651명이라 사전 이력이 제대로 계산된다. 축이 100\% 관측되고 SD 0.37로 분산도
충분하다. \textbf{그런데 라벨과의 상관이 $-$0.011이다.} 축으로 켜면 평균 순위
상관이 $+$0.3478에서 $+$0.3447로 미세하게 내려가고 짝지은 구간이
{[}$-$0.032, $+$0.008{]}로 0을 포함한다. 게임과 같은 처리로 \textbf{보류했다}.
여섯 도메인의 같은 축을 나란히 놓으니 스펙트럼이 드러난다 --- 팝업 $+$0.322,
아이돌 $+$0.227, 게임 $+$0.093, 웹툰 $+$0.087, 도서 $+$0.030, 펀딩 $-$0.011.
\textbf{``만든 이의 이력이 성과를 예측하는가''는 축의 문제가 아니라 시장의
성질이다.} 아이돌은 소속사가 강하게 작동하고 크라우드펀딩은 신인과 베테랑이
구분되지 않는다. 축을 잘못 잰 것이 아니라 그 시장이 그런 것이다. 한편 웹툰이
900건이 되며 자기 순위 상관 $+$0.378로 두 번째로 높은 도메인이 됐고, 전체는
2,181건 서른 셀에 평균 $+$0.3447이다.
\end{abstract}
]

\section{미결 하나를 처리했다}

노트 40에서 펀딩을 넣을 때 매장 노출도를 비워 두었다. 창작자 사전 프로젝트
수로 잡았는데 표본 400건 안에서 세니 평균 0.01 · SD 0.06으로 상수였다. 구조가
다르기 때문이다.

\begin{quote}\fontsize{9.5}{11.5}\selectfont
출판사와 소속사는 소수 주체가 다작하므로 표본 안에서도 사전 건수가 쌓인다.
창작자는 7만 명이 대부분 1건이라 표본 147명 중 2건 이상이 12명뿐이었다.
\end{quote}

목록 3,876쪽을 전부 훑었다.

\begin{table}[t]
\centering\tsize
\caption{텀블벅 전체 색인.}
\begin{tabular}{lr}
\toprule
창작자 & 39,587명 \\
프로젝트 & 76,045건 \\
2건 이상 창작자 & 10,651명 \\
\midrule
펀딩 매장 노출도 관측률 & 100\% \\
축 값 평균 · SD & 0.45 · 0.37 \\
\bottomrule
\end{tabular}
\end{table}

축으로서는 완벽하다. 100\% 관측되고 분산도 충분하다.

\section{그런데 신호가 영이다}

\begin{table}[t]
\centering\tsize
\caption{펀딩 매장 노출도를 켰을 때.}
\begin{tabular}{lr}
\toprule
라벨과의 상관(탈추세) & $-$0.011 \\
\midrule
평균 $\rho$ (끔) & $+$0.3478 \\
평균 $\rho$ (켬) & $+$0.3447 \\
차이 & $-$0.0031 \\
짝지은 95\% 구간 & $[-0.032,\ +0.008]$ \\
P($>$0) & 0.215 \\
\midrule
판정 & \textbf{보류} \\
\bottomrule
\end{tabular}
\end{table}

정확히 계산했는데 정확히 0이다. 게임 매장 노출도와 같은 처리로 꺼 둔다
(노트 28에서 껐다가 노트 37에서 켰다가 노트 43에서 다시 껐다).

\section{여섯 도메인을 나란히 놓으면}

\begin{figure}[t]
\centering
\includegraphics[width=\textwidth]{figs/venue.pdf}
\caption{같은 축, 여섯 시장. 초록이 축으로 쓰는 것, 붉은 것이 꺼 둔 것이다.}
\end{figure}

\begin{table}[t]
\centering\tsize
\caption{매장 노출도 축의 신호.}
\begin{tabular}{lrl}
\toprule
도메인 & r & 무엇으로 쟀나 \\
\midrule
팝업 & \textbf{$+$0.322} & 사람 태깅 \\
아이돌 & $+$0.227 & 소속사 사전 데뷔 \\
게임 & $+$0.093 & 퍼블리셔 사전작 \\
웹툰 & $+$0.087 & 주 연재 횟수 · 작가 이력 \\
도서 & $+$0.030 & 출판사 사전 출간 \\
펀딩 & \textbf{$-$0.011} & 창작자 사전 프로젝트 \\
\bottomrule
\end{tabular}
\end{table}

\claimx{``만든 이의 이력이 성과를 예측하는가''는 축을 잘 쟀느냐의 문제가 아니라
시장의 성질이다. 아이돌은 소속사가 강하게 작동하고($+$0.227) 크라우드펀딩은
신인과 베테랑이 구분되지 않는다($-$0.011).}{펀딩에서 다른 방식으로 이력을 재
(성공률, 누적 모금액, 카테고리 내 순위) 신호가 나오면 축을 잘못 잰 것이다.
지금은 건수만 셌다.}

이 해석이 그럴듯한 이유가 있다. 크라우드펀딩은 진입 장벽이 없고 프로젝트마다
독립적으로 홍보한다 --- 창작자가 앞서 성공했다고 다음 프로젝트가 유리할 구조가
약하다. 반대로 아이돌은 소속사가 방송·유통·팬덤 채널을 이미 갖고 있어 그대로
넘어간다. 도서(출판사)가 $+$0.030으로 낮은 것도 같은 방향에서 읽힌다 ---
출판사 브랜드로 책을 사는 독자가 드물다.

\paragraph{팝업이 또 이겼다}
사람이 기획서를 읽고 매긴 매장 노출도가 $+$0.322로 여섯 중 가장 강하다. 노트
49에서 위키백과가 사람 태깅에 못 미친 것과 같은 방향이다. 다만 여기서는 비교가
공정하지 않다 --- 팝업의 값은 사람이 라벨을 보지 않고 매긴 것이지만, 다른
도메인은 기계적으로 셀 수 있는 것만 셌다.

\section{현재 상태}

\begin{figure}[t]
\centering
\includegraphics[width=\columnwidth]{figs/state.pdf}
\caption{여섯 도메인 서른 셀.}
\end{figure}

\begin{table}[t]
\centering\tsize
\caption{여섯 도메인 2,181건.}
\begin{tabular}{lrrrr}
\toprule
도메인 & n & 자기 $\rho$ & 대상으로 & 출처로 \\
\midrule
아이돌 & 62 & $+$0.472 & $+$0.474 & $+$0.335 \\
\textbf{웹툰} & \textbf{900} & \textbf{$+$0.378} & $+$0.375 & $+$0.314 \\
팝업 & 75 & $+$0.368 & $+$0.449 & $+$0.330 \\
게임 & 482 & $+$0.357 & $+$0.346 & $+$0.342 \\
도서 & 242 & $+$0.202 & $+$0.224 & $+$0.367 \\
펀딩 & 400 & $+$0.190 & $+$0.200 & $+$0.380 \\
\midrule
\multicolumn{2}{l}{평균 $\rho$} & & \multicolumn{2}{r}{$+$0.3447} \\
\multicolumn{2}{l}{자기 $\rho$와의 상관} & & $+$0.956 & 음의 관계 \\
\bottomrule
\end{tabular}
\end{table}

웹툰이 900건이 되며 자기 상관 $+$0.378로 두 번째로 높은 도메인이 됐다. 노트
46에서 처음 넣었을 때 $-$0.125였던 것이 배선 교정 두 번(타깃 폭 $=$ 태그 수,
굿즈 규모 $=$ 작가 수)과 완결작 확보로 여기까지 왔다.

\section{한계}

\begin{itemize}[leftmargin=1.1em,itemsep=0.15em]
\item 펀딩 이력을 \emph{건수}로만 셌다. 성공률이나 누적 모금액으로 재면 다를 수
      있고, 그 데이터는 색인에 이미 들어 있다.
\item 여섯 도메인의 신호를 비교했지만 측정 방식이 제각각이다. 팝업만 사람
      태깅이고 나머지는 기계 집계라 공정한 비교가 아니다.
\item 매장 노출도를 두 도메인에서 꺼 두었는데, 그러면 그 도메인은 공통 축이
      둘이라 쌍별 정렬에서 자유도가 줄어든다. 그 대가를 따로 재지 않았다.
\item 시장 성질이라는 해석은 사후 설명이다. 사전 예측으로 쓴 적이 없다.
\item 웹툰 900건도 전체 완결작의 일부다.
\end{itemize}

\section{다음}

\begin{enumerate}[leftmargin=1.3em,itemsep=0.15em]
\item 펀딩 이력을 성공률·누적 모금액으로 다시 잰다 --- 색인에 데이터가 있다.
\item 매장 노출도를 끈 대가 측정 --- 공통 축이 둘인 도메인의 손해.
\item 큰 도메인을 늘려 구간 반폭이 같은 폭으로 줄어드는지(노트 49의 반증 조건).
\item 시장 성질 해석을 사전 예측으로 쓸 일곱째 도메인 --- 이력이 통할지
      미리 말하고 확인한다.
\end{enumerate}

\vspace{0.6em}
\hrule height 0.3pt
\vspace{0.5em}
{\footnotesize
\textbf{참고문헌}\\[0.3em]
Salganik, M. J., Dodds, P. S., Watts, D. J. (2006). Experimental study of
inequality and unpredictability in an artificial cultural market.
\emph{Science}, 311(5762), 854--856.\\[0.15em]
Merton, R. K. (1968). The Matthew effect in science. \emph{Science},
159(3810), 56--63.\\[0.15em]
Efron, B., Tibshirani, R. J. (1993). \emph{An Introduction to the Bootstrap}.
Chapman \& Hall.\\[0.15em]
Kaufman, S., Rosset, S., Perlich, C. (2012). Leakage in data mining.
\emph{ACM TKDD}, 6(4), 1--21.\\[0.15em]
Ben-David, S. et al. (2010). A theory of learning from different domains.
\emph{Machine Learning}, 79(1--2), 151--175.
}

\end{document}
